\documentclass[10pt]{article}

\usepackage[letterpaper,margin=0.62in]{geometry}
\usepackage{amsmath,amssymb,booktabs,graphicx,microtype,tabularx}
\usepackage[table]{xcolor}
\usepackage{fancyhdr}
\usepackage[hidelinks]{hyperref}

\definecolor{navy}{HTML}{17365D}
\definecolor{brick}{HTML}{A33A2B}
\definecolor{green}{HTML}{236C5B}
\definecolor{purple}{HTML}{4E2A84}
\definecolor{amber}{HTML}{FFF2CC}
\definecolor{pale}{HTML}{F2F4F7}
\definecolor{muted}{HTML}{5F6670}

\setlength{\parindent}{0pt}
\setlength{\parskip}{4pt}
\setlength{\tabcolsep}{5pt}
\renewcommand{\arraystretch}{1.12}

\pagestyle{fancy}
\fancyhf{}
\lhead{\small\color{navy}\textbf{Paper V scalar-stability diagnostic}}
\rhead{\small\color{muted}working evidence}
\cfoot{\small\thepage}
\setlength{\headheight}{14pt}

\newcommand{\statusbanner}{%
  \colorbox{amber}{%
    \parbox{\dimexpr\textwidth-2\fboxsep\relax}{%
      \centering\small\textbf{WORKING DIAGNOSTIC --- NOT PROMOTED PAPER EVIDENCE}%
    }%
  }%
}
\newcommand{\findingbox}[2]{%
  \colorbox{pale}{%
    \parbox{\dimexpr\linewidth-2\fboxsep\relax}{%
      \textcolor{navy}{\textbf{#1}}\quad #2%
    }%
  }%
}
\newcommand{\sourcepath}[1]{\texttt{\detokenize{#1}}}

\input{generated_results.tex}

\begin{document}

\statusbanner
\begin{center}
  \vspace{0.35em}
  {\LARGE\bfseries Regularized equal-time electron--phonon dynamics}\\[0.2em]
  {\Large Scalar instability reproduction and diagnosis}\\[0.35em]
  {\small Paper V working report \(\mid\) generated \ReportGeneratedAt
  \(\mid\) commit \ReportGitCommit}
\end{center}

\findingbox{Outcome.}{The 13-variable scalar Fan--Migdal system reproducibly
crosses \(\max_i|x_i|=10^4\) at
\(t=\BaselineFailureTime\,t_{\mathrm{hop}}^{-1}\) for
\(\lambda=1.5\), \(\gamma=0.5\), and \(v/t_{\mathrm{hop}}=1\).
However, electronic positivity is already lost at
\(t=\BaselinePositivityLoss\,t_{\mathrm{hop}}^{-1}\). Exact-ground-state
contractions plus the Eq.~(112) residual subtraction remain bounded and
    basic-PSD through \(t=400\) in the 13-scalar projection. After adding the two
    real Eq.~(14c) coordinates, however, the relative-mode bosonic uncertainty
    condition is violated at early time. The complete 31-dimensional invariant
    closure confirms the physicality failure. At its first regularized boundary
    crossing, Eq.~(14b)'s correlation source supplies the net outward flux.
    Its causal history is a near-cancellation between the Eq.~(112)
    correlation subtraction and Eq.~(14d)'s bare Pauli-blocking source.}

\findingbox{Correction update.}{A structured full-matrix barrier preserves the
sampled physical cones through \(t=20\); its energy-neutral realization also
retains the post-pulse energy stability of Eq.~(112).}

\vspace{0.25em}
\begin{minipage}[t]{0.49\textwidth}
\textbf{What is supported}
\begin{itemize}
  \item The scalar failure is not a coarse-RK4 artifact.
  \item The Hartree--Fock/zero-correlation state leaves the physical electronic
        domain long before its amplitudes diverge.
  \item The 13 retained scalar derivatives match the direct matrix projection
        of Eqs.~(14a), (14b), (14d), and (14e).
  \item Matrix Eq.~(14c) generates an omitted anomalous two-phonon field.
  \item A two-real-coordinate Eq.~(14c) extension matches the direct matrix
        projection to floating-point precision and delays the HF amplitude
        threshold beyond \(t=360\).
  \item A 31-real-coordinate invariant closure reproduces every generated
        matrix direction and exposes full two-mode bosonic positivity loss.
  \item The first regularized bosonic boundary crossing is driven outward by
        the Eq.~(14b) correlation source.
  \item The boundary-driving correlation is a \(98.9956\%\) causal
        cancellation dominated by Eq.~(112) and the opposing bare
        Pauli-blocking source.
  \item The energy-neutral structured barrier preserves the sampled fermionic
        and bosonic cones through \(t=20\).
\end{itemize}
\end{minipage}\hfill
\begin{minipage}[t]{0.49\textwidth}
\textbf{What is not yet supported}
\begin{itemize}
  \item Equivalence of either the 13D or 15D scalar projection to the complete
        Eqs.~(14a)--(14e); correlation trace/mode-sum sectors remain omitted.
  \item A global minimum-energy root satisfying all higher-moment
        representability constraints.
  \item Long-time and parameter-robust validation of the new
        positivity-preserving correction.
  \item A conclusion that Eq.~(14d) alone causes the failure.
  \item A high-\(U\) conclusion from evidence that currently varies
        electron--phonon coupling \(\lambda\).
  \item Chaos or a strange attractor; the reproduced trajectory is unbounded.
\end{itemize}
\end{minipage}

\begin{figure}[h!]
  \centering
  \includegraphics[width=0.84\linewidth]{trajectory_max_abs.pdf}
  \caption{Maximum scalar-state magnitude for the pinned strong-coupling
  reproduction and weak-coupling control. The threshold is a declared
  diagnostic boundary, not a physical singularity definition.}
\end{figure}

\newpage
\textbf{Immediate decision.}
Treat Eq.~(112) as the amplitude-stabilizing background of a corrected
physical flow. The scalar closure, first boundary flux, Eq.~(14d) causal
history, and a short-time energy-neutral cone correction are resolved for the
pinned protocol. The next classical task is long-time and parameter-robust
validation against exact dimer observables, followed by a correlation-sector
correction that acts on the diagnosed memory channel.

\section*{1.\ Scalar model and deterministic feedback loop}

The scalar-first program has two nested levels. The five-variable Ehrenfest
system uses
\[
x_{\mathrm E}
=
\left(
\Delta n,\,
\operatorname{Re}\rho,\,
\operatorname{Im}\rho,\,
\Delta\operatorname{Re}B,\,
\Delta\operatorname{Im}B
\right),
\]
with the documented invariant
\[
1=\Delta n^2
+4\left[
(\operatorname{Re}\rho)^2+(\operatorname{Im}\rho)^2
\right].
\]
For \(\lambda>1\), the stationary branch changes to
\[
\Delta n_\ast=\pm\sqrt{1-\lambda^{-2}},
\qquad
\operatorname{Re}\rho_\ast=\frac{1}{2\lambda},
\qquad
\Delta\operatorname{Re}B_\ast
=-\frac{2g}{\omega_{\mathrm{ph}}}\Delta n_\ast.
\]

The divergence harness propagates the 13-variable extension in scalar
Eqs.~(87)--(99), adding phonon population/coherence and six real
electron--phonon-correlation coordinates. The pinned initial state is the
Hartree--Fock electronic density,
\(\Delta n=0\), \(\operatorname{Re}\rho=1/2\), with every other coordinate
zero.

\subsection*{Initial residual}

This state has
\[
\lVert F(x_0,0)\rVert_2=\BaselineResidualNorm,
\]
with only two nonzero components:
\begin{center}
\begin{tabular}{@{}lr@{}}
\toprule
Component & Initial derivative \\
\midrule
\ResidualRows
\bottomrule
\end{tabular}
\end{center}
The instantaneous Jacobian has
\[
\max\operatorname{Re}\sigma[J(x_0,0)]
=\BaselineInitialMaxReal.
\]
Because \(F(x_0,0)\ne0\), this number is not a fixed-point stability
certificate. It only shows that the failure is not explained by an immediately
positive eigenvalue at the uncorrelated starting coordinate.

\subsection*{Red-capable command}

\small
\begin{verbatim}
cd paper_5
PYTHONPATH=src python3 -m paper5.stability.reproduce \
  --lambda-ep 1.5 --gamma 0.5 --drive 1.0 \
  --time-step 0.01 --final-time 140 --expect-bounded
\end{verbatim}
\normalsize
The command exits nonzero because the boundedness expectation fails; the first
threshold-crossing component is \BaselineFailureComponent.

\findingbox{Earlier failure.}{The electronic 1-RDM reaches a zero minimum
eigenvalue at \(t=\BaselinePositivityLoss\), approximately \(115\) time units
before the amplitude threshold. The late blow-up therefore occurs after the
trajectory has already escaped the basic physical density-matrix domain.}

\section*{2.\ Solver convergence and growth of local instability}

\begin{center}
\begin{tabular}{@{}lll@{}}
\toprule
Integrator & Step or relative/absolute tolerance & Failure time \\
\midrule
\IntegratorRows
\bottomrule
\end{tabular}
\end{center}

The three RK4 step sizes and three adaptive solvers agree at the scale relevant
to the declared threshold. DOP853 and Radau agree to the displayed seventh
decimal; BDF differs by approximately \(1.4\times10^{-4}\). This strongly
disfavors fixed-step RK4 instability as the sole cause of the scalar failure.

\begin{figure}[h!]
  \centering
  \includegraphics[width=0.89\linewidth]{jacobian_growth.pdf}
  \caption{Maximum real part of the instantaneous scalar Jacobian spectrum
  along the baseline trajectory. These are local eigenvalues, not Lyapunov
  exponents.}
\end{figure}

The largest instantaneous real part is approximately \(1.39\) at \(t=128\),
\(5.19\) at \(t=130\), and \(14.39\) at \(t=130.4\). The trajectory therefore
enters a region of rapidly increasing local amplification before the declared
failure.

\findingbox{Interpretation.}{The ODE solver is tracking a rapidly amplifying
solution of the transcribed scalar equations. This does not yet distinguish a
true finite-time singularity from escape out of the physically admissible
closure manifold.}

\subsection*{Weak-coupling control}

At the same \(\gamma=0.5\), pulse, integration interval, and initial-state
convention, the \(\lambda=0.5\) trajectory remains bounded through \(t=140\),
with \(\max_i|x_i|=\WeakMaximum\). The failure is therefore
coupling-regime-dependent within the scalar model.

\newpage
\section*{3.\ Source ablations and near-failure term balance}

\begin{minipage}[t]{0.51\textwidth}
\centering
\begin{tabular}{@{}lrrrr@{}}
\toprule
Case & \(a_{95}\) & \(a_{97}\) & \(\lVert F_0\rVert\) & \(t_{\mathrm f}\) \\
\midrule
\AblationRows
\bottomrule
\end{tabular}
\end{minipage}\hfill
\begin{minipage}[t]{0.46\textwidth}
\centering
\includegraphics[width=\linewidth]{ablation_failure_times.pdf}
\end{minipage}

Here \(a_{95}\) and \(a_{97}\) scale the initially active density-blocking and
population-source products in scalar Eqs.~(95) and (97). These are diagnostic
ablations, not candidate physical equations. Their retained-coordinate mapping
to the two Pauli-blocking products in matrix Eq.~(14d) is now covered by the
matrix--scalar parity test in Sec.~4.

Deleting either product advances the failure. Deleting both makes
\(F(x_0,0)=0\), yet the strong pulse still triggers failure at \(t=59.92\).
Thus the nonzero initial residual is not the sole cause, and the physical pair
provides a material stabilizing balance.

\subsection*{Large right-hand-side terms near \(t=130\)}

\begin{center}
\small
\begin{tabular}{@{}lr@{}}
\toprule
Contribution & Value \\
\midrule
\CriticalRows
\bottomrule
\end{tabular}
\end{center}

The largest contributions span several equations and include
coherent-displacement/correlation products, electronic-density products, and
phonon-population/coherence products. No single addend dominates every failing
coordinate.

\findingbox{Leading diagnosis.}{The late scalar blow-up is a network-level
nonlinear feedback instability with important cancellation between terms, but
it is downstream of an earlier physical-domain escape. Neither the initial
residual nor either tested source product explains the complete failure chain
alone.}

\newpage
\section*{4.\ Matrix--scalar parity, Eq.~(14c), and invariant closure}

The direct two-site embedding used for the audit is
\[
\rho=
\begin{pmatrix}
(1+\Delta n)/2 & \rho_{\mathrm R}+i\rho_{\mathrm I}\\
\rho_{\mathrm R}-i\rho_{\mathrm I} & (1-\Delta n)/2
\end{pmatrix},
\quad
B=(\Delta B/2,-\Delta B/2),
\quad
\delta\rho_{\mathrm{ph}}=
\begin{pmatrix}n_{\mathrm{ph}}&\rho_{\mathrm{ph}}\\
\rho_{\mathrm{ph}}&n_{\mathrm{ph}}\end{pmatrix}.
\]
For the electron--phonon correlation matrices,
\(C^{(2)}=-C^{(1)}\), \(\operatorname{tr}C^{(1)}=0\), and the real and
imaginary sums/differences of \(C^{(1)}_{12}\) and \(C^{(1)}_{21}\) are the
six coordinates in Eqs.~(102)--(107).

A direct loop implementation of primary Eqs.~(14a)--(14e), with the documented
two-spin factors for the dimer, was evaluated independently of the simplified
scalar formulas. Across \(\ParitySamples\) seeded, basic-admissible random
states, the largest retained-component difference was
\[
\max_{\mathrm{sample},i}
\left|F^{\mathrm{scalar}}_i-
\Pi_iF^{\mathrm{matrix}}\right|
=\ParityMaximumError.
\]

\findingbox{Parity result.}{The signs, conjugations, local Holstein couplings,
spin factors, and normalization of all 13 retained derivatives agree to
floating-point precision. A transcription error inside those retained
coordinates is therefore strongly disfavored.}

\subsection*{Why this is not yet the complete scalarization}

The 13-variable model embeds the anomalous two-phonon fluctuation matrix
\(\delta\bar\rho\) as zero. Matrix Eq.~(14c) does not preserve that condition.
After a first-order \(10^{-3}\) step from the pinned initial state, the norm of
the omitted Eq.~(14c) derivative is already
\(\EqFourteenCNormal\). Thus the model is an exact projection of the retained
equations, but its 13-dimensional slice is not an invariant manifold of the
complete EOM.

\subsection*{Minimal Eq.~(14c) projection}

Dimer symmetry reduces the dynamically generated anomalous relative-mode
field to one complex scalar,
\[
m_-=\delta\bar\rho_{00}-\delta\bar\rho_{01}
    =m_{\mathrm R}+i m_{\mathrm I},
\qquad
\delta\bar\rho=\frac{m_-}{2}
\begin{pmatrix}1&-1\\-1&1\end{pmatrix}.
\]
Its two real equations and its new feedback into the retained correlation
sector are
\[
\dot m_{\mathrm R}=2\omega_{\mathrm{ph}}m_{\mathrm I}
+8g\,\delta_{\mathrm{Im}},
\qquad
\dot m_{\mathrm I}=-2\omega_{\mathrm{ph}}m_{\mathrm R}
-8g\,\delta_{\mathrm{Re}},
\]
\[
\Delta\dot{\mathrm{Im}}^{+}\mapsto
\Delta\dot{\mathrm{Im}}^{+}+2g\,\operatorname{Im}\rho\,m_{\mathrm I},
\qquad
\Delta\dot{\mathrm{Re}}^{+}\mapsto
\Delta\dot{\mathrm{Re}}^{+}+2g\,\operatorname{Im}\rho\,m_{\mathrm R}.
\]
Across the same 100-state audit, all 15 projected derivatives agree with the
matrix RHS to \(\ExtendedParityMaximumError\).

\findingbox{Closure boundary.}{The two coordinates make the Eq.~(14c) sector
tangent to the projected manifold, but do not close the full matrix system.
The pre-existing assumptions \(C^{(2)}=-C^{(1)}\) and
\(\operatorname{tr}C^{(1)}=0\) have a maximum normal derivative
\(\ExtendedPreexistingNormal\) in the random-state audit. The 15D model is
therefore a source-faithful projection, not a complete scalarization.}

\begin{center}
\small
\begin{tabular}{@{}lrrrr@{}}
\toprule
15D protocol & Amplitude \(t_{\mathrm f}\) & Electronic PSD &
Normal-phonon PSD & Boson uncertainty \\
\midrule
\ExtendedRows
\bottomrule
\end{tabular}
\end{center}

The one-mode second-moment condition is
\[
n_-(n_-+1)-|m_-|^2\ge0,
\qquad n_-=n_{\mathrm{ph}}-\rho_{\mathrm{ph}}.
\]
For HF/zero correlations, adding Eq.~(14c) moves the pinned amplitude threshold
from \(130.47\) to \(\ExtendedBaselineFailure\), but the uncertainty condition
fails first. Exact contractions alone still diverge. Exact contractions plus
Eq.~(112) remain amplitude-bounded through \(t=400\), with
\(\max|x_i|=\ExtendedRegularizedMaximum\), but cross the uncertainty boundary
at \(t=\ExtendedRegularizedUncertaintyLoss\) and reach margin
\(\ExtendedRegularizedUncertaintyMinimum\).

\begin{figure}[h!]
  \centering
  \includegraphics[width=0.89\linewidth]{eq14c_protocols.pdf}
  \caption{Matched strong-pulse amplitude diagnostics after adding the two
  Eq.~(14c) coordinates. The late threshold is delayed, not eliminated, for
  HF/zero correlations.}
\end{figure}

\begin{figure}[h!]
  \centering
  \includegraphics[width=0.89\linewidth]{eq14c_uncertainty.pdf}
  \caption{Relative-mode bosonic uncertainty margin for the 15D protocols.
  Negative values cannot represent physical one-mode second moments.}
\end{figure}

\subsection*{Complete 31-scalar invariant closure}

Iterative SVD closure starts from the 15D tangent space, evaluates normal
matrix derivatives at random polynomial probes, and adds every independent
generated direction. It converges from the ambient
\(\ClosedAmbientDimension\)-real-coordinate matrix representation to
\(\ClosedScalarDimension\) coordinates, with disjoint-sample validation
residual \(\ClosedValidationResidual\). The explicit count is
\[
\underbrace{3}_{\rho=\rho^\dagger,\ {\rm tr}\rho=1}
+\underbrace{4}_{B\in\mathbb C^2}
+\underbrace{4}_{N=N^\dagger}
+\underbrace{6}_{A=A^T}
+\underbrace{14}_{C^{(0)},C^{(1)},\
{\rm tr}C^{(0)}={\rm tr}C^{(1)}}
=31.
\]
Here \(N\) is the normal phonon density and \(A\) the anomalous density. No
relative/center or zero-trace correlation component is discarded.

The full two-mode second-moment condition is
\[
\mathcal M_{\rm B}=
\begin{pmatrix}
N^T&A^*\\
A&I+N
\end{pmatrix}\succeq0.
\]
This includes both normal-phonon positivity and anomalous uncertainty in one
test.

\begin{center}
\small
\begin{tabular}{@{}lrrrr@{}}
\toprule
31D protocol & Amplitude \(t_{\mathrm f}\) & Electronic PSD &
Normal-phonon PSD & Full boson moment \\
\midrule
\ClosedRows
\bottomrule
\end{tabular}
\end{center}

The complete HF/zero-correlation closure reaches the amplitude threshold near
\(t=54.46\), but \(\mathcal M_{\rm B}\) loses positivity near \(t=1.33\).
Exact contractions alone reach the amplitude threshold near \(t=141.41\).
Exact contractions plus Eq.~(112) remain bounded through \(t=400\), with
\(\max|x_i|=\ClosedRegularizedMaximum\) and electronic minimum
\(\ClosedRegularizedElectronMinimum\), but the normal-phonon and full
boson-moment minima reach \(\ClosedRegularizedPhononMinimum\) and
\(\ClosedRegularizedBosonMinimum\). The full moment matrix crosses zero at
\(t=\ClosedRegularizedBosonLoss\).

\begin{figure}[h!]
  \centering
  \includegraphics[width=0.89\linewidth]{closed_scalar_protocols.pdf}
  \caption{Amplitude diagnostics for the invariant 31-scalar closure. Only
  exact contractions plus Eq.~(112) remain amplitude-bounded through
  \(t=400\).}
\end{figure}

\begin{figure}[h!]
  \centering
  \includegraphics[width=0.89\linewidth]{closed_scalar_physicality.pdf}
  \caption{Early physicality diagnostics for the residual-subtracted 31D
  closure. Electronic positivity survives, while normal and full bosonic
  positivity fail near \(t=1.49\).}
\end{figure}

\findingbox{Complete-closure result.}{The 31-scalar and complete matrix
trajectories agree. The amplitude instability changes with initialization, but
every unregularized strong protocol exits a physical cone first. Eq.~(112)
removes amplitude divergence and preserves the electronic 1-RDM through
\(t=400\), yet does not preserve the bosonic moment cone.}

The earlier full-matrix probes are recovered exactly: HF/zero reaches the
threshold at \(t=\FullMatrixBaselineFailure\), while exact contractions plus
residual subtraction have electronic minimum
\(\FullMatrixRegularizedElectronMinimum\) and normal-phonon minimum
\(\FullMatrixRegularizedPhononMinimum\) through \(t=140\).

\newpage
\section*{5.\ Bosonic boundary-flux decomposition}

For the 31D exact-contraction plus Eq.~(112) protocol, the first zero of
\(\mathcal M_{\rm B}\) occurs at
\(t_\partial=\BoundaryFluxTime\,t_{\mathrm{hop}}^{-1}\). At this point
\(\lambda_{\min}=\BoundaryFluxMinimum\) and the gap to the next eigenvalue is
\(\BoundaryFluxGap\), so the crossing eigenvalue is simple. If
\(\mathcal M_{\rm B}v=0\), first-order eigenvalue perturbation theory gives
\[
\dot\lambda_{\min}
=v^\dagger\dot{\mathcal M}_{\rm B}v.
\]

\begin{center}
\small
\begin{tabular}{@{}lr@{}}
\toprule
Projected contribution & \(v^\dagger\dot{\mathcal M}_{\rm B}^{(k)}v\) \\
\midrule
Eq.~(14b), correlation source & \(\BoundaryFluxEqFourteenB\) \\
Eq.~(14c), correlation source & \(+\BoundaryFluxEqFourteenC\) \\
Eq.~(14d), direct & \(\BoundaryFluxEqFourteenD\) \\
Eq.~(112), residual subtraction & \(+\BoundaryFluxEqOneTwelve\) \\
\midrule
Total & \(\BoundaryFluxTotal\) \\
\bottomrule
\end{tabular}
\end{center}

The two correlation branches in Eq.~(14b) each contribute
\(\BoundaryFluxEqFourteenBMinus\) and \(\BoundaryFluxEqFourteenBPlus\).
The free Eq.~(14b) and Eq.~(14c) rotations project to zero at the reported
precision. Direct term summation reconstructs the full matrix derivative with
error \(\BoundaryFluxReconstructionError\). A centered finite difference gives
\(\BoundaryFluxFiniteDifference\), with discrepancy
\(\BoundaryFluxFiniteDifferenceError\).

\begin{figure}[h!]
  \centering
  \includegraphics[width=0.82\linewidth]{boson_boundary_flux.pdf}
  \caption{Signed equation-group contributions to
  \(v^\dagger\dot{\mathcal M}_{\rm B}v\) at the first full-boson boundary
  crossing of the residual-subtracted 31D protocol. Negative values point out
  of the positive-semidefinite moment cone.}
\end{figure}

\findingbox{Boundary-flux result.}{The instantaneous outward flux is supplied
by the electron--phonon-correlation source in Eq.~(14b). Eq.~(14c) and the
Eq.~(112) subtraction weakly oppose the crossing. Eq.~(14d) has no direct
first-derivative contribution because \(\mathcal M_{\rm B}\) depends only on
\(N\) and \(A\); it acts indirectly by creating the correlation field entering
Eqs.~(14b) and (14c).}

\subsection*{Causal history through Eq.~(14d)}

Along the realized trajectory, prescribe the other moment sectors and write
Eq.~(14d) as
\[
\dot C=L(t)C+\sum_k S_k(t).
\]
The initial correlation and each independent source are propagated through
the same time-ordered homogeneous evolution generated by \(L(t)\). Their sum
reconstructs the boundary correlation with error
\(\EqFourteenDHistoryReconstruction\); applying the linear Eq.~(14b) flux
functional reconstructs its boundary value with error
\(\EqFourteenDHistoryFluxReconstruction\).

\begin{center}
\small
\begin{tabularx}{0.88\linewidth}{@{}Xr@{}}
\toprule
Causal correlation history & Eq.~(14b) boundary flux \\
\midrule
Eq.~(112), correlation-sector subtraction &
  \(\EqFourteenDHistoryEqOneTwelve\) \\
Eq.~(14d), bare Pauli-blocking source &
  \(+\EqFourteenDHistoryBarePauli\) \\
Eq.~(14d), first anomalous source &
  \(\EqFourteenDHistoryAnomalousOne\) \\
Eq.~(14d), second anomalous source &
  \(+\EqFourteenDHistoryAnomalousTwo\) \\
Eq.~(14d), normal particle source &
  \(+\EqFourteenDHistoryNormalParticle\) \\
Eq.~(14d), normal hole source &
  \(\EqFourteenDHistoryNormalHole\) \\
Propagated initial correlation &
  \(+\EqFourteenDHistoryInitial\) \\
\midrule
Dominant-pair signed sum & \(\EqFourteenDHistoryPairNet\) \\
\bottomrule
\end{tabularx}
\end{center}

\begin{figure}[h!]
  \centering
  \includegraphics[width=0.90\linewidth]{eq14d_history_flux.pdf}
  \caption{Variation-of-constants attribution of the correlation field
  entering the Eq.~(14b) boundary flux. The upper panel exposes the dominant
  cancellation; the lower panel rescales the smaller paired histories. These
  are accumulated causal histories, not direct contributions to
  \(\dot{\mathcal M}_{\rm B}\).}
\end{figure}

The absolute history contributions cancel by
\(\EqFourteenDHistoryCancellationPercent\). The anomalous and
normal-occupation pairs each cancel to approximately \(10^{-8}\); thus the two
paired nonlinear sources are active but do not supply the net outward flux.
The initial Eq.~(112) residual has correlation-sector norm
\(\EqFourteenDHistoryResidualNorm\). Removing only that sector from the
subtraction delays the first crossing from \(\BoundaryFluxTime\) to
\(\EqFourteenDAblatedCrossing\), a delay of
\(\EqFourteenDAblationDelay\,t_{\mathrm{hop}}^{-1}\).

\findingbox{Causal diagnosis.}{For this pinned protocol, the large constant
Eq.~(112) correlation subtraction produces the dominant outward correlation
history, while Eq.~(14d)'s bare Pauli-blocking source nearly cancels it.
Eq.~(112) remains successful as an amplitude regularizer, but applying it
unchanged to the correlation sector advances the bosonic-cone crossing.}

\section*{6.\ Initial conditions, positivity, and residual subtraction}

\begin{center}
\small
\begin{tabular}{@{}lrrrr@{}}
\toprule
Initial-state protocol & \(\lVert F_0\rVert_2\) &
Amplitude \(t_{\mathrm f}\) & Electronic PSD loss & \(\max|x_i|\) \\
\midrule
\InitializationRows
\bottomrule
\end{tabular}
\end{center}

Here \(>400\) means no event through the requested final time. ``Electronic
PSD loss'' is the first zero of the smaller eigenvalue of the \(2\times2\)
electronic 1-RDM. The corresponding retained phonon check uses the eigenvalues
\(n_{\mathrm{ph}}\pm|\rho_{\mathrm{ph}}|\).

\begin{figure}[h!]
  \centering
  \includegraphics[width=0.89\linewidth]{initialization_protocols.pdf}
  \caption{Matched strong-pulse scalar protocols. Correlated initialization
  changes the late-time behavior, while exact contractions plus Eq.~(112)
  remain bounded through \(t=400\).}
\end{figure}

\subsection*{Source-connected stationary root}

Exact ground-state contractions at phonon cutoff 16 seed a deterministic
least-squares solve of \(F(x_\ast,0)=0\). The selected root has residual
\(\StationaryResidual\), basic electronic and phonon positivity, and energy
\(\StationaryEnergy\), compared with exact truncated ground-state energy
\(\ExactSeedEnergy\). It is bounded through \(t=400\), but loses electronic
positivity at \(t=30.56\). Because higher-moment representability constraints
are not yet parameterized, this is a source-connected physical candidate, not
a proof of the global minimum over every algebraic root.

\subsection*{Eq.~(112) diagnostic regularization}

The working source proposes
\[
\dot x(t)=F(x(t),t)-F(x_0,0),
\]
with \(x_0\) obtained from exact ground-state contractions. This makes that
contracted initialization stationary by construction. For the strong
\(\lambda=1.5\) scalar protocol, the maximum magnitude through \(t=400\) is
\(\RegularizedMaximum\); the minimum electronic and retained phonon
eigenvalues are respectively \(\RegularizedElectronMinimum\) and
\(\RegularizedPhononMinimum\).
This 13D check omits \(m_-\), so it does not test the bosonic uncertainty
condition exposed by the 15D extension.

\begin{figure}[h!]
  \centering
  \includegraphics[width=0.89\linewidth]{regularized_physicality.pdf}
  \caption{Basic positivity diagnostics for the residual-subtracted strong
  scalar protocol. Both displayed minima remain above zero through \(t=400\).}
\end{figure}

\findingbox{Isolation result.}{The \(t\approx130.47\) blow-up is not an
irreducible numerical singularity. In the scalar projection it can be removed
by a source-documented correlated initialization plus residual subtraction.
The unresolved problem is upgrading this successful amplitude correction into
a representability-preserving evolution of the now-complete 31D closure.}

\section*{7.\ Structured cone correction}

The direct correction uses the ten real coordinates of a Hermitian
\(\Delta\dot N\) and complex-symmetric \(\Delta\dot A\).  Their structured
lift \(\mathcal S(y)\) obeys the same block identity as the bosonic moment
matrix.  With barrier rate \(\beta=5\), safety margin
\(h_\star=10^{-5}\), and \(\kappa=0\), each right-hand-side evaluation solves
\[
\begin{aligned}
\underset{y\in\mathbb R^{10}}{\operatorname{minimize}}\quad&
\frac12\lVert y\rVert_2^2,\\
\operatorname{subject\ to}\quad&
D+\mathcal S(y)
+\beta\!\left(\mathcal M_{\rm B}-h_\star I\right)
-\kappa I\succeq0.
\end{aligned}
\]
The matrix inequality constrains the complete active eigenspace.  A
single-minimum-eigenvector rule reaches a near degeneracy of the two lowest
bosonic eigenvalues near \(t=3.60\), where eigenvector switching produces
solver chattering.  Constraint generation over violated eigenvectors removes
that coordinate singularity and supplies a minimum-norm point for the
accumulated half-spaces.

For a correction confined to \(N\) and \(A\), Eq.~(22) receives the
instantaneous contribution
\[
\Delta\dot E
=\omega_{\rm ph}
\left(\Delta\dot N_{00}+\Delta\dot N_{11}\right).
\]
The energy-neutral chart therefore imposes
\(\Delta\dot N_{00}+\Delta\dot N_{11}=0\) by optimizing in the null space of
this row.

\begin{center}
\small
\begin{tabular}{@{}lrrrr@{}}
\toprule
Protocol & \(\min\lambda(\mathcal M_{\rm B})\) &
\(\min\lambda(\rho)\) & \(\max|x_i|\) &
\(\max_{t\geq4}|E(t)-E(4)|\)\\
\midrule
\ConeCorrectionRows
\bottomrule
\end{tabular}
\end{center}

\begin{figure}[h!]
  \centering
  \includegraphics[width=0.91\linewidth]{cone_correction_audit.pdf}
  \caption{Pinned 31D correction audit through \(t=20\).  The direct and
  energy-neutral matrix barriers preserve the bosonic cone.  The middle panel
  exposes the energy cost of the unconstrained correction and the recovery
  obtained by the trace-neutral moment velocity.}
\end{figure}

\findingbox{Short-time correction result.}{The unconstrained direct barrier is
active on \(\ConeDirectActiveFraction\) of sampled times and reaches correction
norm \(\ConeDirectMaximumCorrection\), but its post-pulse energy drift is
\(\ConeDirectEnergyDrift\).  The energy-neutral barrier is active on
\(\ConeEnergyActiveFraction\) of samples, reaches norm
\(\ConeEnergyMaximumCorrection\), has maximum direct energy injection
\(\ConeEnergyFlux\), and reduces post-pulse drift to
\(\ConeEnergyDrift\).  All sampled energy-neutral subproblems converged
(\(\ConeEnergyNonconverged\) unresolved samples).}

\section*{8.\ Diagnosis boundaries and next calculations}

\begin{enumerate}
  \item \textbf{Extend the energy-neutral run.} Repeat the \(t=20\) result
  through the amplitude-validation horizon and across step, tolerance,
  cutoff, pulse, and coupling refinements.

  \item \textbf{Parameterize representability.} Enforce electronic
  \(0\preceq\rho\preceq I\), the full bosonic covariance uncertainty
  condition---including \(n_-(n_-+1)\ge|m_-|^2\)---and compatible
  electron--phonon moment constraints during both root selection and time
  propagation.

  \item \textbf{Separate correction mechanisms.} Compare residual subtraction,
  the energy-neutral moment barrier, and a regularized Eq.~(14d) correction
  one at a time against exact dimer dynamics.

  \item \textbf{Repeat across \(\lambda\).} The strong scalar case is
  encouraging, but the same correction has small late-time positivity
  violations in the weak control and the complete matrix probe. It is not yet
  a universal prescription.

  \item \textbf{Only then test chaos.} A strange attractor requires bounded
  attraction, a robust positive largest Lyapunov exponent, volume contraction,
  and fractal geometry on a validated physical flow.
\end{enumerate}

\subsection*{Closed-form opportunity}

Equations (14b), (14c), and (14e) have forced-oscillator integral forms, and
Eq.~(14d) has a time-ordered propagator form when the other sectors are
prescribed. The most useful analytic target is not a general elementary
trajectory; it is the scalar fixed-point system, its exact Jacobian, and the
stability boundaries of each physical branch:
\[
F(x_\ast,0)=0,
\qquad
J_\ast=\left.\frac{\partial F}{\partial x}\right|_{x_\ast},
\qquad
\delta x(t)=e^{J_\ast(t-t_0)}\delta x(t_0).
\]

\subsection*{Evidence boundary}

The explicit local reproduction varies electron--phonon coupling
\(\lambda=2g^2/(t_{\mathrm{hop}}\omega_{\mathrm{ph}})\). It must not be
reported as a high-\(U\) result until a separate Hubbard interaction \(U\) is
defined in the scalar model and varied under a matched protocol.

\section*{9.\ Reproducibility and parameter manifest}

\begin{tabularx}{\textwidth}{@{}lX@{}}
\toprule
Field & Recorded value \\
\midrule
Artifact status & Working diagnostic; not promoted paper evidence \\
Generated & \ReportGeneratedAt \\
Git commit & \ReportGitCommit \\
Baseline & \(t_{\mathrm{hop}}=1\), \(\gamma=0.5\),
\(\lambda=1.5\), \(v/t_{\mathrm{hop}}=1\),
\(\Delta t=0.01\), \(t_{\mathrm f}=140\) \\
Initial state & Hartree--Fock electronic density,
\(\Delta n=0\), \(\operatorname{Re}\rho=1/2\), all remaining scalar
coordinates zero \\
Failure definition & First time \(\max_i|x_i|>10^4\) or a nonfinite scalar
component appears \\
Weak control & Same protocol with \(\lambda=0.5\) \\
Scalar source & \sourcepath{/Users/jakestrobel/Downloads/Dynamics_on_the_Hubbard_DIMER.pdf},
Eqs.~(78)--(82), (84)--(107), and (112) \\
Supporting derivation & \sourcepath{/Users/jakestrobel/Downloads/Electron_phonon_interactions___chiral_phonons.pdf} \\
Primary matrix source & G. Riva, J. Simoni, and Y. Ping,
\href{https://arxiv.org/abs/2606.22233}{arXiv:2606.22233},
Eqs.~(14a)--(14e) \\
Matrix-to-scalar status & Retained 13D projection verified; 15D Eq.~(14c)
projection verified but incomplete; invariant 31D coordinate system matches
the complete matrix flow \\
\bottomrule
\end{tabularx}

\subsection*{Short source hashes}

\begin{tabular}{@{}ll@{}}
\toprule
Source & SHA-256 prefix \\
\midrule
Scalar implementation & \CodeHashShort \\
Working note & \NoteHashShort \\
Hubbard-dimer PDF & \HubbardPdfHashShort \\
Chiral-phonon PDF & \ChiralPdfHashShort \\
\bottomrule
\end{tabular}

\subsection*{Stable paths}

\small
\begin{tabularx}{\textwidth}{@{}lX@{}}
\toprule
Role & Path \\
\midrule
LaTeX source &
\sourcepath{MATH/paper_facing/paper_V_high_u_gkba/paper_v_scalar_stability_diagnostic.tex} \\
Numerical generator &
\sourcepath{paper_5/src/paper5/stability/report.py} \\
Scalar model &
\sourcepath{paper_5/src/paper5/stability/hubbard_dimer.py} \\
Matrix reference &
\sourcepath{paper_5/src/paper5/stability/matrix_reference.py} \\
Exact contractions &
\sourcepath{paper_5/src/paper5/stability/exact_reference.py} \\
Initial-state and Eq.~(112) workflow &
\sourcepath{paper_5/src/paper5/stability/initial_conditions.py} \\
Final PDF &
\sourcepath{output/pdf/paper_v_scalar_stability_diagnostic.pdf} \\
Full manifest &
\sourcepath{output/pdf/paper_v_scalar_stability_diagnostic_manifest.json} \\
\bottomrule
\end{tabularx}
\normalsize

\subsection*{Build and validation commands}

\small
\begin{verbatim}
cd paper_5
python3 -m pytest -q
PYTHONPATH=src python3 -m paper5.stability.report
\end{verbatim}
\normalsize

\findingbox{Append policy.}{Add new numerical results through the scalar
harness and report generator, preserve prior completed evidence in the
manifest, and keep hypotheses visibly separate from confirmed diagnostics.
Manuscript adoption remains a user decision.}

\end{document}
